%%% One-side print:
\documentclass[11pt,a4paper]{report}
\usepackage[top=25mm,bottom=25mm,right=25mm,left=30mm,head=12.5mm,foot=12.5mm]{geometry}
\let\openright=\clearpage

%%% Two-side print:
%\documentclass[11pt,a4paper,twoside,openright]{report}
%\usepackage[top=25mm,bottom=25mm,right=25mm,left=30mm,head=12.5mm,foot=12.5mm]{geometry}
%\let\openright=\cleardoublepage

%%% Definition of some useful macros
%%% Tento soubor obsahuje definice různých užitečných maker a prostředí %%%
%%% Další makra připisujte sem, ať nepřekáží v ostatních souborech.     %%%
%%% This file contains definitions of various useful macros and environments      %%%
%%% Assign additional macros here so that they do not interfere with other files. %%%

\usepackage[a-2u]{pdfx}     % výsledné PDF bude ve standardu PDF/A-2u
                            % resulting PDF will be in the PDF / A-2u standard

\usepackage{ifpdf}
\usepackage{ifxetex}
\usepackage{ifluatex}

%%% Nastavení pro použití samostatné bibliografické databáze.
%%% Settings for using a separate bibliographic database.
\usepackage[
   backend=biber
 ,style=iso-authoryear
  % ,style=iso-numeric
  ,sortlocale=cs_CZ
  ,alldates=iso
  ,bibencoding=UTF8
  ,maxnames=2
  ,maxbibnames=99
  ,block=ragged
]{biblatex}
\let\cite\parencite
\renewcommand*{\multinamedelim}{, \addspace}
\renewcommand*{\finalnamedelim}{\addspace a \addspace}

\bibliography{bibliography}

%% Přepneme na českou sazbu, fonty Latin Modern a kódování češtiny
\ifthenelse{\boolean{xetex}\OR\boolean{luatex}}
   { % use fontspec and OpenType fonts with utf8 engines
			\usepackage[english,slovak,czech]{babel}
			\usepackage[autostyle,english=british,czech=quotes]{csquotes}
			\usepackage{fontspec}
			\defaultfontfeatures{Ligatures=TeX,Scale=MatchLowercase}
   }
   {
			\usepackage[english,slovak,czech]{babel}
			\usepackage{lmodern}
			\usepackage[T1]{fontenc}
			\usepackage{textcomp}
			\usepackage[utf8]{inputenc}
			\usepackage[autostyle,english=british,czech=quotes]{csquotes}
	 }
\ifluatex
\makeatletter
\let\pdfstrcmp\pdf@strcmp
\makeatother
\fi

%%% Další užitečné balíčky (jsou součástí běžných distribucí LaTeXu)
\usepackage{amsmath}        % rozšíření pro sazbu matematiky / extension for math typesetting
\usepackage{amsfonts}       % matematické fonty / mathematical fonts
\usepackage{amssymb}        % symboly / symbols
\usepackage{amsthm}         % sazba vět, definic apod. / typesetting of sentences, definitions, etc.
\usepackage{bm}             % tučné symboly (příkaz \bm) / bold symbols (\bm command)
\usepackage{graphicx}       % vkládání obrázků / graphics inserting
\usepackage{listings}       % vylepšené prostředí pro strojové písmo / improved environment for source codes typesetting
\usepackage{fancyhdr}       % prostředí pohodlnější nastavení hlavy a paty stránek / environment for more comfortable adjustment of the head and foot of the pages
\usepackage{icomma}         % inteligetní čárka v matematickém módu / intelligent comma in math mode
\usepackage{dcolumn}        % lepší zarovnání sloupců v tabulkách / better alignment of columns in tables
\usepackage{booktabs}       % lepší vodorovné linky v tabulkách / better horizontal lines in tables
\makeatletter
\@ifpackageloaded{xcolor}{
   \@ifpackagewith{xcolor}{usenames}{}{\PassOptionsToPackage{usenames}{xcolor}}
  }{\usepackage[usenames]{xcolor}} % barevná sazba / color typesetting
\makeatother
\usepackage{multicol}       % práce s více sloupci na stránce / work with multiple columns on a page
\usepackage{caption}
\usepackage{enumitem}
\setlist[itemize]{noitemsep, topsep=0pt, partopsep=0pt}
\setlist[enumerate]{noitemsep, topsep=0pt, partopsep=0pt}
\setlist[description]{noitemsep, topsep=0pt, partopsep=0pt}

\usepackage{tocloft}
\setlength\cftparskip{0pt}
\setlength\cftbeforechapskip{1.5ex}
\setlength\cftfigindent{0pt}
\setlength\cfttabindent{0pt}
\setlength\cftbeforeloftitleskip{0pt}
\setlength\cftbeforelottitleskip{0pt}
\setlength\cftbeforetoctitleskip{0pt}
\renewcommand{\cftlottitlefont}{\Huge\bfseries\sffamily}
\renewcommand{\cftloftitlefont}{\Huge\bfseries\sffamily}
\renewcommand{\cfttoctitlefont}{\Huge\bfseries\sffamily}

% vyznaceni odstavcu
% differentiation of new paragraphs
\parindent=0pt
\parskip=11pt

% zakaz vdov a sirotku - jednoradkovych pocatku ci koncu odstavcu na prechodu mezi strankami
% Prohibition of widows and orphans - single-line beginning and end of paragraph at the transition between pages
\clubpenalty=1000
\widowpenalty=1000
\displaywidowpenalty=1000

% nastaveni radkovani
% setting of line spacing
\renewcommand{\baselinestretch}{1.20}

% nastaveni pro nadpisy - tucne a bezpatkove
% settings for headings - bold and sans serif
\usepackage{sectsty}    
\allsectionsfont{\sffamily}

% nastavení hlavy a paty stránek
% page head and foot settings
\makeatletter
\if@twoside%
    \fancypagestyle{fancyx}{%
			\fancyhf{}                                     
      \fancyhead[RE]{\rightmark}                  
      \fancyhead[LO]{\leftmark}                  
      \fancyfoot[RO,LE]{\thepage}                    
      \renewcommand{\headrulewidth}{.5pt}            
      \renewcommand{\footrulewidth}{.5pt}            
    }
    \fancypagestyle{plain}{%
			\fancyhf{}                                     
    	\fancyfoot[RO,LE]{\thepage}                    
    	\renewcommand{\headrulewidth}{0pt}             
    	\renewcommand{\footrulewidth}{0.5pt}
    }          
\else
    \fancypagestyle{fancyx}{%
			\fancyhf{}                                     
      \fancyhead[R]{\leftmark}                  
      \fancyfoot[R]{\thepage}                    
      \renewcommand{\headrulewidth}{.5pt}            
      \renewcommand{\footrulewidth}{.5pt}            
    }
    \fancypagestyle{plain}{%                       
    	\fancyhf{} % clear all header and footer fields
    	\fancyfoot[R]{\thepage}                    
    	\renewcommand{\headrulewidth}{0pt}             
    	\renewcommand{\footrulewidth}{0.5pt}
    }          
\fi
\renewcommand*{\cleardoublepage}{\clearpage\if@twoside \ifodd\c@page\else
	\hbox{}%
	\thispagestyle{empty}%
	\newpage%
	\if@twocolumn\hbox{}\newpage\fi\fi\fi
}
\makeatother

% Tato makra přesvědčují mírně ošklivým trikem LaTeX, aby hlavičky kapitol
% sázel příčetněji a nevynechával nad nimi spoustu místa. Směle ignorujte.
% These macros convince with a slightly ugly LaTeX trick to make chapter headers
% bet more sane and didn't miss a lot of space above them. Be boldly ignore it.
\makeatletter
\def\@makechapterhead#1{
  {\parindent \z@ \raggedright \sffamily
   \Huge\bfseries \thechapter. #1
   \par\nobreak
   \vskip 20\p@
}}
\def\@makeschapterhead#1{
  {\parindent \z@ \raggedright \sffamily
   \Huge\bfseries #1
   \par\nobreak
   \vskip 20\p@
}}
\makeatother

% Trochu volnější nastavení dělení slov, než je default.
% Slightly looser hyphenation setting than default.
\lefthyphenmin=2
\righthyphenmin=2

% Zapne černé "slimáky" na koncích řádků, které přetekly, abychom si jich lépe všimli.
% Turns on the black "snails" at the ends of the lines that overflowed to get us noticed them better.
% \overfullrule=1mm

%% Balíček hyperref, kterým jdou vyrábět klikací odkazy v PDF,
%% ale hlavně ho používáme k uložení metadat do PDF (včetně obsahu).
%% Většinu nastavítek přednastaví balíček pdfx.
%% A hyperref package that can be used to produce clickable links in PDF,
%% but we mainly use it to store metadata in PDF (including content).
%% Most settings are preset by the pdfx package.
\hypersetup{unicode}
\hypersetup{breaklinks=true}
\hypersetup{hidelinks}

% Přejmenování objektů autoref pro odkazování v textu
\renewcommand{\chapterautorefname}{Kapitola}
\renewcommand{\sectionautorefname}{Sekce}
\renewcommand{\subsectionautorefname}{Podsekce}
\renewcommand{\tableautorefname}{Tabulka}
\renewcommand{\figureautorefname}{Obrázek}
\renewcommand{\equationautorefname}{Rovnice}



\renewcommand{\UrlBreaks}{\do\/\do\=\do\+\do\-\do\_\do\ \do\a\do\b\do\c\do\d%
\do\e\do\f\do\g\do\h\do\i\do\j\do\k\do\l\do\m\do\n\do\o\do\p\do\q\do\r\do\s%
\do\t\do\u\do\v\do\w\do\x\do\y\do\z\do\A\do\B\do\C\do\D\do\E\do\F\do\G\do\H%
\do\I\do\J\do\K\do\L\do\M\do\N\do\O\do\P\do\Q\do\R\do\S\do\T\do\U\do\V\do\W%
\do\X\do\Y\do\Z\do\1\do\2\do\3\do\4\do\5\do\6\do\7\do\8\do\9\do\0}

\renewcommand{\mid}{|}

%%% Prostředí pro sazbu kódu, případně vstupu/výstupu počítačových
%%% programů. (Vyžaduje balíček listings -- fancy verbatim.)
%%% Environment for source code typesetting, or computer input/output
%%% programs. (Requires package listings - fancy verbatim.)
\lstnewenvironment{code}{\lstset{basicstyle=\small, frame=single}}{}

%%% User-defined balíčky
\usepackage{dsfont}

\theoremstyle{definition}
\newtheorem{definition}{Definice}[section]
\newcommand{\definitionautorefname}{Definice}

\usepackage{float}
\usepackage{tikz}
\usetikzlibrary{fit,arrows.meta,automata,positioning,shapes.geometric,trees}

\usepackage{tikzscale}

\usepackage{forest}

\usepackage[ruled]{algorithm2e}
\usepackage{subfig}
\usepackage{pdflscape}
\usepackage{minted}

\usepackage{neuralnetwork}
\newcommand{\xin}[2]{$x_#2$}
\newcommand{\xout}[2]{$\hat x_#2$}

\DeclareMathOperator{\tr}{\text{tr}}
\SetAlgorithmName{Algoritmus}{Algoritmus}




%%% DEFINITION OF BASIC VARIABLES
\def\TypPrace{BP}                % bakalářská práce/bachelor thesis
%\def\TypPrace{DP}               % diplomová práce/master thesis
\def\Jazyk{cze}                  % čeština/czech
%\def\Jazyk{slo}                 % slovenština/slovak
%\def\Jazyk{eng}                 % angličtina/english

%%% Název práce v jazyce práce (přesně podle zadání)
%%% Title of the thesis in the language used in the text (exact according to assignment)
\def\NazevPrace{Variační autoenkodér a úlohy pozorování v latentním prostoru}

%%% Tento soubor obsahuje definice různých užitečných maker a prostředí %%%
%%% Další makra připisujte sem, ať nepřekáží v ostatních souborech.     %%%
%%% This file contains definitions of various useful macros and environments      %%%
%%% Assign additional macros here so that they do not interfere with other files. %%%

\usepackage[a-2u]{pdfx}     % výsledné PDF bude ve standardu PDF/A-2u
                            % resulting PDF will be in the PDF / A-2u standard

\usepackage{ifpdf}
\usepackage{ifxetex}
\usepackage{ifluatex}

%%% Nastavení pro použití samostatné bibliografické databáze.
%%% Settings for using a separate bibliographic database.
\usepackage[
   backend=biber
 ,style=iso-authoryear
  % ,style=iso-numeric
  ,sortlocale=cs_CZ
  ,alldates=iso
  ,bibencoding=UTF8
  ,maxnames=2
  ,maxbibnames=99
  ,block=ragged
]{biblatex}
\let\cite\parencite
\renewcommand*{\multinamedelim}{, \addspace}
\renewcommand*{\finalnamedelim}{\addspace a \addspace}

\bibliography{bibliography}

%% Přepneme na českou sazbu, fonty Latin Modern a kódování češtiny
\ifthenelse{\boolean{xetex}\OR\boolean{luatex}}
   { % use fontspec and OpenType fonts with utf8 engines
			\usepackage[english,slovak,czech]{babel}
			\usepackage[autostyle,english=british,czech=quotes]{csquotes}
			\usepackage{fontspec}
			\defaultfontfeatures{Ligatures=TeX,Scale=MatchLowercase}
   }
   {
			\usepackage[english,slovak,czech]{babel}
			\usepackage{lmodern}
			\usepackage[T1]{fontenc}
			\usepackage{textcomp}
			\usepackage[utf8]{inputenc}
			\usepackage[autostyle,english=british,czech=quotes]{csquotes}
	 }
\ifluatex
\makeatletter
\let\pdfstrcmp\pdf@strcmp
\makeatother
\fi

%%% Další užitečné balíčky (jsou součástí běžných distribucí LaTeXu)
\usepackage{amsmath}        % rozšíření pro sazbu matematiky / extension for math typesetting
\usepackage{amsfonts}       % matematické fonty / mathematical fonts
\usepackage{amssymb}        % symboly / symbols
\usepackage{amsthm}         % sazba vět, definic apod. / typesetting of sentences, definitions, etc.
\usepackage{bm}             % tučné symboly (příkaz \bm) / bold symbols (\bm command)
\usepackage{graphicx}       % vkládání obrázků / graphics inserting
\usepackage{listings}       % vylepšené prostředí pro strojové písmo / improved environment for source codes typesetting
\usepackage{fancyhdr}       % prostředí pohodlnější nastavení hlavy a paty stránek / environment for more comfortable adjustment of the head and foot of the pages
\usepackage{icomma}         % inteligetní čárka v matematickém módu / intelligent comma in math mode
\usepackage{dcolumn}        % lepší zarovnání sloupců v tabulkách / better alignment of columns in tables
\usepackage{booktabs}       % lepší vodorovné linky v tabulkách / better horizontal lines in tables
\makeatletter
\@ifpackageloaded{xcolor}{
   \@ifpackagewith{xcolor}{usenames}{}{\PassOptionsToPackage{usenames}{xcolor}}
  }{\usepackage[usenames]{xcolor}} % barevná sazba / color typesetting
\makeatother
\usepackage{multicol}       % práce s více sloupci na stránce / work with multiple columns on a page
\usepackage{caption}
\usepackage{enumitem}
\setlist[itemize]{noitemsep, topsep=0pt, partopsep=0pt}
\setlist[enumerate]{noitemsep, topsep=0pt, partopsep=0pt}
\setlist[description]{noitemsep, topsep=0pt, partopsep=0pt}

\usepackage{tocloft}
\setlength\cftparskip{0pt}
\setlength\cftbeforechapskip{1.5ex}
\setlength\cftfigindent{0pt}
\setlength\cfttabindent{0pt}
\setlength\cftbeforeloftitleskip{0pt}
\setlength\cftbeforelottitleskip{0pt}
\setlength\cftbeforetoctitleskip{0pt}
\renewcommand{\cftlottitlefont}{\Huge\bfseries\sffamily}
\renewcommand{\cftloftitlefont}{\Huge\bfseries\sffamily}
\renewcommand{\cfttoctitlefont}{\Huge\bfseries\sffamily}

% vyznaceni odstavcu
% differentiation of new paragraphs
\parindent=0pt
\parskip=11pt

% zakaz vdov a sirotku - jednoradkovych pocatku ci koncu odstavcu na prechodu mezi strankami
% Prohibition of widows and orphans - single-line beginning and end of paragraph at the transition between pages
\clubpenalty=1000
\widowpenalty=1000
\displaywidowpenalty=1000

% nastaveni radkovani
% setting of line spacing
\renewcommand{\baselinestretch}{1.20}

% nastaveni pro nadpisy - tucne a bezpatkove
% settings for headings - bold and sans serif
\usepackage{sectsty}    
\allsectionsfont{\sffamily}

% nastavení hlavy a paty stránek
% page head and foot settings
\makeatletter
\if@twoside%
    \fancypagestyle{fancyx}{%
			\fancyhf{}                                     
      \fancyhead[RE]{\rightmark}                  
      \fancyhead[LO]{\leftmark}                  
      \fancyfoot[RO,LE]{\thepage}                    
      \renewcommand{\headrulewidth}{.5pt}            
      \renewcommand{\footrulewidth}{.5pt}            
    }
    \fancypagestyle{plain}{%
			\fancyhf{}                                     
    	\fancyfoot[RO,LE]{\thepage}                    
    	\renewcommand{\headrulewidth}{0pt}             
    	\renewcommand{\footrulewidth}{0.5pt}
    }          
\else
    \fancypagestyle{fancyx}{%
			\fancyhf{}                                     
      \fancyhead[R]{\leftmark}                  
      \fancyfoot[R]{\thepage}                    
      \renewcommand{\headrulewidth}{.5pt}            
      \renewcommand{\footrulewidth}{.5pt}            
    }
    \fancypagestyle{plain}{%                       
    	\fancyhf{} % clear all header and footer fields
    	\fancyfoot[R]{\thepage}                    
    	\renewcommand{\headrulewidth}{0pt}             
    	\renewcommand{\footrulewidth}{0.5pt}
    }          
\fi
\renewcommand*{\cleardoublepage}{\clearpage\if@twoside \ifodd\c@page\else
	\hbox{}%
	\thispagestyle{empty}%
	\newpage%
	\if@twocolumn\hbox{}\newpage\fi\fi\fi
}
\makeatother

% Tato makra přesvědčují mírně ošklivým trikem LaTeX, aby hlavičky kapitol
% sázel příčetněji a nevynechával nad nimi spoustu místa. Směle ignorujte.
% These macros convince with a slightly ugly LaTeX trick to make chapter headers
% bet more sane and didn't miss a lot of space above them. Be boldly ignore it.
\makeatletter
\def\@makechapterhead#1{
  {\parindent \z@ \raggedright \sffamily
   \Huge\bfseries \thechapter. #1
   \par\nobreak
   \vskip 20\p@
}}
\def\@makeschapterhead#1{
  {\parindent \z@ \raggedright \sffamily
   \Huge\bfseries #1
   \par\nobreak
   \vskip 20\p@
}}
\makeatother

% Trochu volnější nastavení dělení slov, než je default.
% Slightly looser hyphenation setting than default.
\lefthyphenmin=2
\righthyphenmin=2

% Zapne černé "slimáky" na koncích řádků, které přetekly, abychom si jich lépe všimli.
% Turns on the black "snails" at the ends of the lines that overflowed to get us noticed them better.
% \overfullrule=1mm

%% Balíček hyperref, kterým jdou vyrábět klikací odkazy v PDF,
%% ale hlavně ho používáme k uložení metadat do PDF (včetně obsahu).
%% Většinu nastavítek přednastaví balíček pdfx.
%% A hyperref package that can be used to produce clickable links in PDF,
%% but we mainly use it to store metadata in PDF (including content).
%% Most settings are preset by the pdfx package.
\hypersetup{unicode}
\hypersetup{breaklinks=true}
\hypersetup{hidelinks}

% Přejmenování objektů autoref pro odkazování v textu
\renewcommand{\chapterautorefname}{Kapitola}
\renewcommand{\sectionautorefname}{Sekce}
\renewcommand{\subsectionautorefname}{Podsekce}
\renewcommand{\tableautorefname}{Tabulka}
\renewcommand{\figureautorefname}{Obrázek}
\renewcommand{\equationautorefname}{Rovnice}



\renewcommand{\UrlBreaks}{\do\/\do\=\do\+\do\-\do\_\do\ \do\a\do\b\do\c\do\d%
\do\e\do\f\do\g\do\h\do\i\do\j\do\k\do\l\do\m\do\n\do\o\do\p\do\q\do\r\do\s%
\do\t\do\u\do\v\do\w\do\x\do\y\do\z\do\A\do\B\do\C\do\D\do\E\do\F\do\G\do\H%
\do\I\do\J\do\K\do\L\do\M\do\N\do\O\do\P\do\Q\do\R\do\S\do\T\do\U\do\V\do\W%
\do\X\do\Y\do\Z\do\1\do\2\do\3\do\4\do\5\do\6\do\7\do\8\do\9\do\0}

\renewcommand{\mid}{|}

%%% Prostředí pro sazbu kódu, případně vstupu/výstupu počítačových
%%% programů. (Vyžaduje balíček listings -- fancy verbatim.)
%%% Environment for source code typesetting, or computer input/output
%%% programs. (Requires package listings - fancy verbatim.)
\lstnewenvironment{code}{\lstset{basicstyle=\small, frame=single}}{}

%%% User-defined balíčky
\usepackage{dsfont}

\theoremstyle{definition}
\newtheorem{definition}{Definice}[section]
\newcommand{\definitionautorefname}{Definice}

\usepackage{float}
\usepackage{tikz}
\usetikzlibrary{fit,arrows.meta,automata,positioning,shapes.geometric,trees}

\usepackage{tikzscale}

\usepackage{forest}

\usepackage[ruled]{algorithm2e}
\usepackage{subfig}
\usepackage{pdflscape}
\usepackage{minted}

\usepackage{neuralnetwork}
\newcommand{\xin}[2]{$x_#2$}
\newcommand{\xout}[2]{$\hat x_#2$}

\DeclareMathOperator{\tr}{\text{tr}}
\SetAlgorithmName{Algoritmus}{Algoritmus}




%%% Jméno autora
%%% Author's name - First name Surname
\def\AutorPrace{Tomáš Faltejsek}

%%% Rok odevzdání a měsíc (slovně)
%%% Year of submission and month (verbally) - month YYYY
\def\DatumOdevzdani{květen 2023}

%%% Vedoucí práce: Jméno a příjmení s~tituly
%%% Supervisor: First name and surname with titles
\def\Vedouci{Ing. Ondřej Vadinský, Ph.D.}

%%% Konzultant práce: Jméno a příjmení s~tituly
%%% Consultant: First name and surname with titles
\def\Konzultant{full consultant’s name (incl. degrees)}

%%% Studijní program
%%% Study program
\def\StudijniProgram{Aplikovaná informatika}

%%% Studijní program - specializace
%%% Study program - specialization
\def\Specializace{}

%%% Studijní obor
%%% Field of study
\def\StudijniObor{Aplikovaná informatika}

%%% Nepovinné poděkování (vedoucímu práce, konzultantovi, tomu, kdo zapůjčil software, literaturu apod.)
%%% Optional thanks (the supervisor, the consultant, the borrower of software, literature, etc.)
\def\Podekovani{%
Thanks.
}

%%% Abstract (recommended range of about 150-250 words; this is not a thesis assignment)
\def\Abstrakt{%
Jedním z předních rysů lidské inteligence je intuice a schopnost představovat si nové objekty. Variační autoenkodér je inovací na poli pravděpodobnostních modelů, umožňující architekturu modelů schopných syntézy zcela nových dat s využitím pozorování atributů v latentním prostoru. Teoretická charakteristika a možnosti využití variačního autoenkodéru jsou předmětem této bakalářské práce.
}
\def\KlicovaSlova{keyword, important term, another topic, and another one}


%%% Titulní strana a různé povinné informativní strany
%%% Title page and various mandatory information pages
\begin{document}
%%% Titulní strana práce a další povinné informační strany
%%% Title page of the thesis and other obligatory information pages

%%% Titulní strana práce
%%% Title page of the thesis

\pagestyle{empty}
\hypersetup{pageanchor=false}

\begin{center}
\Huge\sffamily
\VSE\\
\FIS

\vspace{\stretch{1}}

\ifstringequal{\Jazyk}{eng}{
	\includegraphics[width=.5\textwidth]{img/FIS_2_logo_2_rgb_EN}
}{
	\includegraphics[width=.5\textwidth]{img/FIS_2_logo_2_rgb_EN}
	% \includegraphics[width=.5\textwidth]{img/FIS_2_logo_2_rgb_CZ} TODO: Add CZ variant of the logo
}

\vspace{\stretch{2}}

\bfseries\NazevPrace

\vspace{8mm}
\mdseries\TypPraceText

\vspace{8mm}
\large
\begin{tabular}{rl}
\StudijniProgramText: & \StudijniProgram \\
\ifthenelse{\equal{\Specializace}{}}{%
	% empty value
	}{
	\rule{0pt}{6mm}%
	\SpecializaceText: & \Specializace \\
}
\ifthenelse{\equal{\StudijniObor}{}}{%
	% empty value
	}{
	\rule{0pt}{6mm}%
	\StudijniOborText: & \StudijniObor \\
}
\end{tabular}

\vspace{\stretch{8}}

\begin{tabular}{rl}
\AutorText: & \AutorPrace \\
\noalign{\vspace{2mm}}
\VedouciText: & \Vedouci \\
\ifthenelse{\equal{\Konzultant}{}}{%
	% empty value
	}{
	\rule{0pt}{6mm}%
	\KonzultantText: & \Konzultant \\
}
\end{tabular}

\vspace{8mm}
\Praha, \DatumOdevzdani
\end{center}


%%% Poděkování
%%% Acknowledgments
\hypersetup{pageanchor=true}
\cleardoublepage
\pagestyle{plain}
\openright
\vspace*{\fill}
\section*{\PodekovaniText}
\noindent
\Podekovani
\vspace{1cm}


%%% Povinná informační strana práce
%%% Obligatory information page of the thesis
\openright
\section*{Abstract}
\noindent
\Abstrakt
\subsection*{Keywords}
\noindent
\KlicovaSlova

\openright


%%% Strana s automaticky generovaným obsahem práce
%%% A page with automatically generated content of the thesis
\setcounter{tocdepth}{2}
\tableofcontents

%%% Seznam obrázků v práci
%%% List of figures in the thesis
\openright
\listoffigures
Note: Add a list of figures if the number of figures in the thesis text exceeds 20. A list of diagrams is applicable only if the author distinguishes between a figure and a diagram. The list of diagrams is included if the number of diagrams exceeds 20. This thesis template does not distinguish between a figure and a diagram.

%%% Seznam tabulek v práci (volitelně)
%%% List of tables in the thesis (optionally)
\clearpage
\listoftables
Note: Add a list of tables if the number of tables used in the thesis exceeds 20. 

%%% Seznam zkratek v práci (volitelně)
%%% List of abbreviations in the thesis (optionally)
\chapter*{\SeznamZkratek}

\begin{multicols}{2}
\raggedright
\begin{description}
\item [BCC] Blind Carbon Copy
\item [CC] Carbon Copy
\item [CERT] Computer Emergency Response Team
\item [CSS] Cascading Styleheets
\item [DOI] Digital Object Identifier
\item [HTML] Hypertext Markup Language
\item [REST] Representational State Transfer
\item [SOAP] Simple Object Access Protocol
\item [URI] Uniform Resource Identifier
\item [URL] Uniform Resource Locator
\item [XML] eXtended Markup Language
\end{description}
\end{multicols}

Note: Add a list of abbreviations if the number of abbreviations used in the thesis exceeds 20 and the abbreviations used are not common.

\pagestyle{fancyx}
%%% Jednotlivé kapitoly práce jsou pro přehlednost uloženy v samostatných souborech
%%% The individual chapters of the thesis are stored in separate files for clarity
{%
\pagestyle{plain}
\chapter*{Introduction}
\addcontentsline{toc}{chapter}{Introduction}

Introduction is a compulsory part of the bachelor's / diploma thesis. The introduction is an introduction to the topic. It elaborates the chosen topic, briefly puts it into context (there may also be a description of the motivation to write the work) and answers the question why the topic was chosen. It puts the topic into context and justifies its necessity and the topicality of the solution. It contains an explicit goal of the work. The text of the thesis goal is identical with the text that is given in the bachelor's thesis assignment, ie with the text that is given in the InSIS system and which is also given in the Abstract section.

Part of the introduction is also a brief introduction to the process of processing the work (a separate part of the actual text of the work is devoted to the method of processing). The introduction may also include a description of the motivation to write the work.

The introduction to the diploma thesis must be more elaborate - this is stated in more detail in the Requirements of the diploma thesis within the Intranet for FIS students.

Here are some sample chapters that recommend how a bachelor's / master's thesis should be set. They primarily describe the use of the \LaTeX{} template, but general advice will also serve users of other systems well.
}
% \include{...}
% \include{...}

\chapter{Východiska variačního autoenkodéru}
% TODO: Move to separate files
\section{Redukce dimenzionality}
\section{Neuronové sítě}
\section{Autoenkodér}
\section{Pravděpodobnostní modely}
\section{Hluboké modely s latentní reprezentací}
\section{Kullback-Lieblerova divergence}

% \include{...}
% \include{...}
\chapter{Variační autoenkodér}
\section{ELBO}
\section{Reparametrizační trik}
\section{Formalizace}
\section{Model neuronové sítě}
\section{Nedostatky a omezení}
\section{Rozšíření a aktuální stav poznání}
\section{Pozorování v latentním prostoru}


\chapter{Úlohy pozorování v latentním prostoru}
\section{Generativní modelování obrazových dat}
\section{Rekonstrukce obrazových dat}
\section{Interpolace vět}
\section{Detekce anomálií}
\section{Syntéza tabulárních dat}
\section{Komprese}

\chapter{Experimenty s modelem variačního autoenkodéru}
\section{Generativní modelování obrazových dat}
\subsection{Vymezení problémové oblasti}
\subsection{Datová sada a předzpracování}
\subsection{Nastavení experimentu}
\subsection{Návrh modelu}
\subsection{Evaluace}
\subsection{Diskuze}
\section{Interpolace vět}

{%
\pagestyle{plain}
\chapter*{Závěr}
\addcontentsline{toc}{chapter}{Závěr}

Tématem práce byl \emph{variační autoenkodér a úlohy pozorování v latentním prostoru}.

V první kapitole práce bylo prezentováno široké spektrum možných aplikací modelu variačního autoenkodéru formou stručného shrnutí závěrů dostupných publikací.
Tato kapitola měla za cíl představit čtenáři způsob, jakým lze variační autoenkodér využít, a motivovat ho, aby pokračoval v četbě teoretické části práce.

\textbf{Teoretický úsek výkladu} je považován za \textbf{hlavní přínos této práce}, jelikož postupně staví a detailně interpretuje teoretické aspekty variačního autoenkodéru pomocí následujících kapitol:
\begin{itemize}
    \item \autoref{chap:prereqs} uvádí seznam východisek variačního autoenkodéru, které jsou v oblasti strojového učení stabilně ukotveny.
    \item \autoref{chap:autoencoder} představuje autoenkodér – architekturu modelu strojového učení, na jehož principech staví variační autoenkodér. Byl popsán základní princip fungování autoenkodéru, jeho jednotlivé typy a způsob jejich regularizace. Kapitola je zakončena výkladem o stochastickém autoenkodéru, který slouží jako předchůdce pro uvedení variačního autoenkodéru. 
    \item \autoref{chap:vae} definuje princip fungování variačního autoenkodéru s důrazem na interpretaci jeho teoretických aspektů. Detailně byla rozebrána rovněž účelová funkce variačního autoenkodéru – konkrétně její dva prvky: KL divergence a chyba rekonstrukce, včetně role, kterou v trénovacím procesu modelu zastávají. Závěrem kapitoly byl krátce shrnut aktuální stav poznání variačního autoenkodéru, jeho existující rozšíření a omezení.
\end{itemize}

V tomto jde práce nad rámec monografie autorů variačního autoenkodéru \textcite{Kingma2019} a propojuje variační autoenkodér s architekturou různých typů autoenkodéru a dalších definic z oblasti strojového učení (autoři monografie předpokládají znalost této látky a tak ji ve své monografii neuvádí).

Následně byl na základě zavedené teorie prakticky implementován \textbf{ilustrační} model variačního autoenkodéru pro generativní úlohu obrazových dat MNIST.
Latentní prostor naučeného modelu byl formou vizualizací analyzován a interpretován. Závěrem kapitoly jsou diskutovány možnosti evaluace generativních modelů variačního autoenkodéru.

Výsledkem práce je \textbf{ucelený} výklad o variačním autoenkodéru, který na jednom místě shrnul možnosti jeho multidisciplinárních aplikací, definoval a interpretoval jeho teoretické aspekty a prakticky implementoval ilustrační model variačního autoenkodéru pro generativní úlohu obrazových dat MNIST.
K implementovanému modelu a jeho vizualizacím jsou rovněž přiloženy kompletní zdrojové kódy pro snadnou replikaci.
}


%%% Seznam použité literatury
%%% Bibliography
%% This applies if a separate bibliographic database is used
\printbibliography[title={\bibname},heading={bibintoc}]

%% This is true when using thebibliography environment
%% The following can be recommended for compiling citation data:
%%     https://knihovna.vse.cz/citace/priklady/
%%     https://www.citace.com/
%\openright
%\phantomsection
%\addcontentsline{toc}{chapter}{\bibname}
%\begin{thebibliography}{99}
%\bibitem{Cermak2018}ČERMÁK, Radim, SMUTNÝ, Zdeněk. A Framework for Cultural Localization of Websites and for Improving Their Commercial Utilization. In:  \emph{Global Observations of the Influence of Culture on Consumer Buying Behavior} [online]. Hershey~: IGI Global, 2018, s. 206--232. ISBN 978-1-5225-2727-5. DOI: 10.4018/978-1-5225-2727-5.
%
%\bibitem{Hladik2018}HLADÍK, Milan, ČERNÝ, Michal. The Shape of the Optimal Value of a Fuzzy Linear Programming Problem. In: \emph{Fuzzy Logic in Intelligent System Design} [online]. Cancum, 16.10.2017 -- 18.10.2017. Cham~: Springer, 2018, s. 281--286. Advances in Intelligent Systems and Computing 648. ISBN 978-3-319-67136-9. DOI: 10.1007/978-3-319-67137-6\_31.
%
%\bibitem{Jasek2018}JAŠEK, Pavel, VRANÁ, Lenka, ŠPERKOVÁ, Lucie, SMUTNÝ, Zdeněk, KOBULSKÝ, Marek. Modeling and Application of Customer Lifetime Value in Online Retail. \emph{Informatics} [online]. 2018, roč. 5, č. 1. 22 s. eISSN 2227-9709. DOI: 10.3390/informatics5010002. Dostupné také z: \url{http://www.mdpi.com/2227-9709/5/1/2/pdf}.
%
%\bibitem{Pecakova2018}PECÁKOVÁ, Iva. \emph{Statistika v terénních průzkumech}. 3. přeprac. vyd. Praha~: Professional Publishing, 2018. 254 s. ISBN 978-80-88260-10-3.
%\end{thebibliography}


%%% Přílohy k práci, existují-li. Každá příloha musí být alespoň jednou
%%% odkazována z vlastního textu práce. Přílohy se číslují.
%%% Attachments to thesis, if any. Each attachment must be referenced at 
%%% least once in your own text. The appendices are numbered.
\part*{\Prilohy\thispagestyle{empty}}
\appendix
\chapter{Zdrojové kódy modelů}
\label{app:vae_model_source_code}

Přiložený soubor \emph{\lstinline|bachelor-thesis-source-code.zip|} obsahuje:

\begin{itemize}
    \item \emph{\lstinline|vae_mnist_train.py|}
    \item \emph{\lstinline|vae_mnist_train.pynb|}
    \item \emph{save}
    \item \emph{model}
    \item \emph{loaders}
\end{itemize}

Pro spuštění trénovací fáze a vizualizací jsou potřeba následující programy a balíčky:
\begin{itemize}
    \item Python 3.12.
    \item TensorFlow 2.12.0.
    \item Matplotlib 3.7.1.
    \item Numpy 1.24.0
    \item Jupyter notebook 6.5.4.
\end{itemize}

Trénovací fáze se spustí souborem \lstinline|vae_mnist_train.py|, vizualizace pak \lstinline|vae_mnist_train.pynb|.
% \include{...}
% \include{...}

\end{document}
